{\setlength{\parskip}{8mm}
{\huge{\textbf{Resumen}}}

{\large{

El \gls{timple} es un pequeño instrumento profundamente arraigado en la cultura de las Islas Canarias. A pesar de su relevancia en el folclore local, este nunca llegó a popularizarse al nivel de otros instrumentos como la guitarra o el ukelele. Esta situación se ve agravada por la escasez de recursos accesibles y gratuitos que faciliten su aprendizaje y práctica.

Este proyecto propone el desarrollo de una aplicación web orientada a la creación, difusión y práctica de canciones para timple. Para ello, la aplicación contará con un conjunto de herramientas para la creación y reproducción de canciones, páginas informátivas con utilidades para introducirse al instrumento y un espacio donde acceder a todos los recursos anteriores.

A lo largo del documento se expone de forma estructurada el proceso de desarrollo del proyecto, abarcando las fases de análisis, diseño, implementación, pruebas y despliegue, así como otros aspectos relevantes como la normativa y legislación, herramientas, metodología y planificación utilizadas.

}}

\newpage

{\huge{\textbf{Abstract}}}


{\large{

The \emph{\gls{timple}} is a small instrument deeply rooted in the culture of the  Canary Islands. Although relevant in local culture, it did not have the chance to be as popular as other instruments such as the guitar or ukulele. The lack of reachable and free resources for the learning and practice of the instrument made the situation even worse.

The aim of this project is to develop of a web application that facilitates the creation, destribution and practice of music for the \emph{timple}. To achieve this, the application counts with a set of tools for the creating and practicing of songs, informative pages with utilities that ease the first steps with the instrument, and an easy to access space to the resources mentioned.

Throughout this document you will follow the development of this project, getting through phases such as analysis, design, implementation, tests and deployment, as other relevant aspects such as normative and legilation, tools, methodology and planificacion used.

}

}}